\section{Descenso de propiedades topológicas}
\label{VIII.4}

\begin{theoreme}
\label{VIII.4.1}
Sean $g\colon Y'\to Y$ un morfismo y $Z$ un subconjunto de~$Y$.
Supongamos que~$g$ es plano y que existe un morfismo
cuasicompacto~$f\colon X\to Y$ tal que~$Z=f(X)$ (N.B.: si $Y$ es
noetheriano, esta última condición está implicada por
\ifthenelse{\boolean{orig}}
{$Z$ es constructible}
{«$Z$ es constructible»}). Entonces se tiene
$$
g^{-1}(\overline{Z})=\overline{g^{-1}(Z)}\quoi.
$$
\end{theoreme}

\ifthenelse{\boolean{orig}}
{Puede suponerse~$Y$ afín y después $Y'$ afín. Como $Y$ es afín,
$X$ es una \emph{reunión
\marginpar{206}
finita de abiertos afines~$X_{i}$ y, reemplazando~$X$ por la suma de
los~$X_{i}$, puede suponerse asimismo que~$X$ es afín. Sean~$A,A',B$ los
anillos de~$Y,Y',X$, sea $B'=B\otimes_{A}A'$ el de
$X'=X\times_{Y}Y'$, sea $I$ el núcleo de~$A\to B$ y sea $I'$ el núcleo
de~$A'\to B'$. Así, los subconjuntos cerrados de~$Y$ y~$Y'$ definidos
por estos ideales son, respectivamente, la adherencia de~$Z=f(X)$ y la
adherencia de~$Z'=f'(X')=g^{-1}(Z)$. Se quiere probar que esta última es
igual a~$g^{-1}(\overline{Z})$, lo que resultará de~$I'=IA'$; a su vez,
esta igualdad es consecuencia de la planitud de~$A'$ sobre~$A$.}}
{Puede suponerse~$Y$ afín y después $Y'$ afín. Como $Y$ es afín,
$X$ es una reunión
\marginpar{206}
finita de abiertos afines~$X_{i}$ y, reemplazando~$X$ por la suma de
los~$X_{i}$, puede suponerse asimismo que~$X$ es afín. Sean~$A,A',B$ los
anillos de~$Y,Y',X$, sea $B'=B\otimes_{A}A'$ el de
$X'=X\times_{Y}Y'$, sea $I$ el núcleo de~$A\to B$ y sea $I'$ el núcleo
de~$A'\to B'$. Así, los subconjuntos cerrados de~$Y$ y~$Y'$ definidos
por estos ideales son, respectivamente, la adherencia de~$Z=f(X)$ y la
adherencia de~$Z'=f'(X')=g^{-1}(Z)$. Se quiere probar que esta última es
igual a~$g^{-1}(\overline{Z})$, lo que resultará de~$I'=IA'$; a su vez,
esta igualdad es consecuencia de la planitud de~$A'$ sobre~$A$.}

\begin{corollaire}
\label{VIII.4.2}
Sean $g\colon Y'\to Y$ un morfismo plano y cuasicompacto y $Z'$ un
subconjunto cerrado de~$Y'$ saturado para la relación de equivalencia
conjuntista definida por~$g$. Entonces
$Z'=g^{-1}(\overline{g(Z')})$.
\end{corollaire}

\ifthenelse{\boolean{orig}}{}
{\enlargethispage{.5cm}}%
En efecto, se tiene $Z'=g^{-1}(Z)$, con~$Z=g(Z')$. Puede aplicarse
entonces~\Ref{VIII.4.1}, observando que la condición impuesta a~$Z$ en
\Ref{VIII.4.1} se satisface tomando como~$X$ el preesquema~$Z'$ provisto de
la estructura reducida inducida por~$Y'$. (El hecho de que~$g$ sea
cuasicompacto asegura entonces que el morfismo inducido
$f\colon Z'\to Y$ sea cuasicompacto).

El enunciado~\Ref{VIII.4.2} significa también que \emph{la topología de
~$g(Y')$ inducida por~$Y$ es la topología cociente de la de~$Y'$}. En
particular:

\begin{corollaire}
\label{VIII.4.3}
Sea $g\colon Y'\to Y$ un morfismo fielmente plano y cuasicompacto. Entonces
$g$ hace de~$Y$ un espacio topológico cociente de~$Y'$; es decir, para un
subconjunto~$Z$ de~$Y$, $Z$ es cerrado (\resp abierto) si y solo si lo es
$Z'=g^{-1}(Z)$.
\end{corollaire}

Recordemos ahora que dos elementos
\ifthenelse{\boolean{orig}}
{$a,b$, de~$Y'$}
{$a,b$ de~$Y'$}
tienen la misma imagen en~$Y$ si y solo si
\ifthenelse{\boolean{orig}}
{son}
{son}
de la forma $p_{1}(c),p_{2}(c)$ para un elemento conveniente~$c$ de
$Y''=Y'\times_{Y}Y'$. Se sigue que, si $g$ es sobreyectivo, se tiene un
diagrama \emph{exacto} de conjuntos
$$
\xymatrix@C=.5cm{\cal{P}(Y)\ar[r] & \cal{P}(Y')\ar@/^2pt/[r]\ar@/_2pt/[r] &
\cal{P}(Y'')}\quoi,
$$
donde, para todo conjunto~$E$, se denota por~$\cal{P}(E)$ el conjunto de
sus subconjuntos. Sentado esto, \Ref{VIII.4.3} puede interpretarse también
del modo siguiente:

\begin{corollaire}[Descenso de los subconjuntos abiertos \resp cerrados]
\label{VIII.4.4}
Sea
\marginpar{207}%
\index{descenso de los subconjuntos abiertos (\resp cerrados)|hyperpage}%
$g\colon Y'{\to} Y$ como en~\Ref{VIII.4.3}. Para todo preesquema~$X$, sea
$\Ouv(X)$ \resp $\Fer(X)$ el conjunto de sus subconjuntos abiertos \resp
el conjunto de sus subconjuntos cerrados. Entonces se tienen diagramas
\emph{exactos} de aplicaciones entre conjuntos (deducidos de~$g$ y de las
dos proyecciones de~$Y''=Y'\times_{Y}Y'$):
\vspace*{-5pt}
\begin{gather*}
\xymatrix@C=.5cm{\Ouv(Y)\ar[r] & \Ouv(Y')\ar@/^2pt/[r]\ar@/_2pt/[r] &
\Ouv(Y'')}\\[-3pt]
\xymatrix@C=.5cm{\Fer(Y)\ar[r] & \Fer(Y')\ar@/^2pt/[r]\ar@/_2pt/[r] &
\Fer(Y'')}
\end{gather*}
\end{corollaire}

Se tiene el complemento siguiente de~\Ref{VIII.4.3}:

\begin{corollaire}
\label{VIII.4.5}
Sea $g\colon Y'\to Y$ como en~\Ref{VIII.4.3}, y sea~$Z$ un subconjunto
de~$Y$ tal que exista un morfismo cuasicompacto~$f\colon X\to Y$ de
imagen~$Z$ (por ejemplo, $Z$ constructible e~$Y$ noetheriano). Para que
$Z$ sea un subconjunto localmente cerrado de~$Y$, es necesario y suficiente
que $Z'=g^{-1}(Z)$ sea un subconjunto localmente cerrado de~$Y'$.
\end{corollaire}

Basta probar la suficiencia. Sea $Y_{1}$ el subpreesquema cerrado de~$Y$
dado por la adherencia de~$Z$ provista de la estructura reducida inducida,
y sea~$Y_{1}'=Y_{1}\times_{Y}Y'$ el subpreesquema cerrado de~$Y'$ imagen
inversa de~$Y_{1}$. Su conjunto subyacente es la imagen inversa
$g^{-1}(Y_{1})=g^{-1}(\overline{Z})$ y, por tanto, es igual, en virtud
de~\Ref{VIII.4.1}, a~$\overline{Z'}$. Como $Z'$ es localmente cerrado
en~$Y'$, es abierto en~$\overline{Z'}$ y, por tanto, abierto en~$Y_{1}'$.
Ahora bien, este último es fielmente plano y cuasicompacto sobre~$Y_{1}$;
por~\Ref{VIII.4.3} se concluye que~$Z$ es abierto en~$Y_{1}$, es decir,
en~$\overline{Z}$, lo que significa que $Z$ es localmente cerrado.

\begin{corollaire}
\label{VIII.4.6}
Sea $g\colon S'\to S$ un morfismo fielmente plano y cuasicompacto,
$f\colon X\to Y$ un~$S$-morfismo y $f'\colon X'\to Y'$ el
$S'$-morfismo que se deduce de él por cambio de base. Supongamos que~$f'$
es una aplicación abierta (\resp una aplicación cerrada, \resp
\ifthenelse{\boolean{orig}}
{cuasicompacta}
{cuasicompacta}
y un homeomorfismo sobre su imagen, \resp un homeomorfismo); entonces~$f$
tiene la misma propiedad.
\end{corollaire}

Como $Y'$ es fielmente plano y cuasicompacto sobre~$Y$, puede suponerse
$Y=S$. Sea~$Q$
\marginpar{208}
un subconjunto de~$X$. Se tiene entonces (denotando por~$h$ el morfismo de
proyección $X'\to X$):
$$
g^{-1}(f(Q))=f'(h^{-1}(Q))\quoi.
$$

Si $Q$ es abierto (\resp cerrado), lo mismo sucede con~$h^{-1}(Q)$ y, por
tanto, también con~$f'(h^{-1}(Q))$ si se supone que~$f'$ es una aplicación
abierta (\resp cerrada); por la fórmula precedente y~\Ref{VIII.4.3}, lo
mismo sucede entonces con~$f(Q)$. Esto demuestra las dos primeras
afirmaciones de~\Ref{VIII.4.6}. Queda por examinar el caso en que~$f'$ es
un homeomorfismo sobre su imagen, y demostrar entonces que~$f$ es un
homeomorfismo sobre su imagen. (El caso de un homeomorfismo resultará
entonces de~\Ref{VIII.3.1}). En virtud
\ifthenelse{\boolean{orig}}
{de~\Ref{VIII.3.1}}
{de~\Ref{VIII.3.2}}
$f$ es inyectivo; queda por probar que la aplicación $X\to f(X)$ es
abierta. Se sabe ya que~$f$ es cuasicompacto en virtud
de~\Ref{VIII.3.3}. Basta probar que, para todo subconjunto cerrado~$Z$
de~$X$, se tiene $Z=f^{-1}(\overline{f(Z)})$; esto equivale a la fórmula
análoga para las imágenes inversas por la aplicación sobreyectiva
$h\colon X'\to X$, es decir, a
$$
\ifthenelse{\boolean{orig}}
{Z'=f'^{-1}(g^{-1}(\overline{f(Z)})}
{Z'=f'^{-1}(g^{-1}(\overline{f(Z)}))}\quoi,
$$
donde se pone $Z'=h^{-1}(Z)$. Ahora bien, en virtud de~\Ref{VIII.4.1}
aplicado al subconjunto~$f(Z)$ de~$Y$, se tiene
$g^{-1}(\overline{f(Z)})=\overline{g^{-1}(f(Z))}$, y la fórmula que hay
que demostrar equivale a
$$
Z'=f'^{-1}(\overline{f'(Z')})\quoi,
$$
la cual resulta de la hipótesis de que~$f'$ es un homeomorfismo sobre su
imagen.

N.B.: en este último razonamiento, suponiendo ya que~$f$ es cuasicompacto,
no se ha utilizado que~$g$ sea cuasicompacto, sino solamente que
\ifthenelse{\boolean{orig}}
{$g'$}
{$g$}
es fielmente plano. Por tanto, bajo esta hipótesis puede descenderse la
propiedad «homeomorfismo sobre su imagen» o «homeomorfismo», o también,
gracias al razonamiento precedente, la propiedad «$f'$ es cuasicompacto y
hace de~$f'(X')$ un espacio topológico cociente de~$X'$».

Diremos que un morfismo~$f\colon X\to Y$ de preesquemas es
\emph{universalmente abierto} (\resp \emph{universalmente cerrado},
\resp \emph{universalmente bicontinuo},
\ifthenelse{\boolean{orig}}{etc...}{etc.})
\index{universalmente abierto (\resp universalmente cerrado, \resp
universalmente bicontinuo, etc.)|hyperpage}%
si, para todo cambio de base~$Y'\to Y$, $f'\colon X'\to Y'$ es un
morfismo abierto (\resp cerrado, \resp un
\ifthenelse{\boolean{orig}}
{homeom.}
{homeomorfismo}
sobre su espacio imagen). Se deduce entonces de~\Ref{VIII.4.6}:

\begin{corollaire}
\label{VIII.4.7}
Bajo
\marginpar{209}
las condiciones de~\Ref{VIII.4.6}, para que~$f$ sea universalmente
abierto (\resp universalmente cerrado, \resp un homeomorfismo universal
sobre su imagen, \resp un homeomorfismo universal), es necesario y
suficiente que~$f'$ lo sea.
\end{corollaire}

\begin{corollaire}
\label{VIII.4.8}
Bajo las condiciones de~\Ref{VIII.4.6}, para que~$f$ sea separado
(\resp propio), es necesario y suficiente que~$f'$ lo sea.
\end{corollaire}

Decir que~$f$ es separado significa que el morfismo diagonal
$X\to X\times_{Y}X$ es cerrado, o también universalmente cerrado, y, por
tanto, la primera afirmación
\ifthenelse{\boolean{orig}}
{}
{de}
~\Ref{VIII.4.8} resulta de~\Ref{VIII.4.7}. Decir que~$f$ es propio
significa que~$f$ satisface las condiciones: a)~$f$ es de tipo finito;
b)~$f$ es separado; c)~$f$ es universalmente cerrado. La condición a)
desciende bien en virtud de~\Ref{VIII.3.3}; la b), también, por lo que
acaba de verse; finalmente, la c), asimismo, por~\Ref{VIII.4.7}.

\begin{remarques}
\label{VIII.4.9}
Recordemos que, cuando $g\colon Y'\to Y$ es un morfismo plano de tipo
finito, con~$Y$ localmente noetheriano, entonces~$g$ es un morfismo
abierto
(\ifthenelse{\boolean{orig}}{VI~6.6}{\Ref{IV.6.6}}),
lo que es un resultado más preciso que~\Ref{VIII.4.3}. Obsérvese, sin
embargo, que, si~$f$ es un morfismo fielmente plano y cuasicompacto de
preesquemas noetherianos, entonces~$f$ no es, en general, un morfismo
abierto. Sea, por ejemplo, $Y$ un esquema irreducible cuyo punto
genérico~$y$ no sea abierto (por ejemplo, una curva algebraica), y tomemos
como~$Y'$ el esquema suma~$Y\amalg\Spec(\kres(y))$. Entonces la imagen por
el morfismo estructural~$Y'\to Y$ del subconjunto
abierto~$\Spec(\kres(y))$ no es un subconjunto abierto de~$Y$. El lector
observará asimismo que diversos enunciados de la presente exposición se
vuelven falsos si se abandona en ellos la hipótesis de que el morfismo
fielmente plano considerado sea también cuasicompacto. El caso típico que
refuta los enunciados es aquel en que se toma como~$Y'$ el esquema suma de
los espectros de los anillos locales de los puntos de~$Y$. Por ejemplo,
tomando de nuevo como~$Y$ una curva algebraica irreducible y como~$Z$ el
subconjunto de~$Y$ reducido al punto genérico, su imagen inversa en~$Y'$ es
abierta sin que~$Z$ sea abierto.
\end{remarques}

\setcounter{subsection}{9}
\subsection{}
\label{VIII.4.10}
Diversos
\marginpar{210}
enunciados de la presente exposición siguen siendo válidos reemplazando la
hipótesis de que~$Y'$ sea plano sobre~$Y$ por la siguiente: existe un
\ifthenelse{\boolean{orig}}
{módulo}
{módulo}
de tipo finito~$F$ sobre~$Y'$, de soporte~$Y'$, plano con respecto a~$Y$;
la hipótesis de planitud fiel se reemplazará entonces por la precedente,
más la hipótesis de que~$Y'\to Y$ sea sobreyectivo. Esto se aplica a las
dos primeras afirmaciones de~\Ref{VIII.1.10}, a~\Ref{VIII.3.3},
\Ref{VIII.3.5}, \Ref{VIII.4.1} y, por consiguiente, a todos los resultados
del presente número.
